\section{Conclusion}
\begin{frame}{Conclusion}
\begin{itemize}
\item We have introduced a basic workflow for obtaining free energies of binding
\item We have taken many short-cuts in the interest of time
\item The principles in this workshop can be transferred to other software packages
\item Relative free energies require a second calculation with a different ligand to be useful!
\end{itemize}
\end{frame}

\begin{frame}{Acknowledgements}
Although this workshop is my own work, It would not have been possible without the help of...
\newline
\begin{columns}
\column{0.5\textwidth}
My supervisory team:
\begin{itemize}
\item Prof. Chris Hayes
\item Prof. Jonathan Hirst
\item Dr. Christof J\"{a}ger
\item Dr. David Rogers
\end{itemize}

My funding group:
\begin{itemize}
\item BBSRC DTP
\item The University of Nottingham
\end{itemize}
\column{0.5\textwidth}

My colleagues:
\begin{itemize}
\item Dr. Bang Huynh
\item Aiden Cranney
\item Comp. Chem. Nottingham! 
\end{itemize}
\begin{multicols}{2}
\centering
\begin{figure}
\includegraphics[height=0.08\textheight]{../Graphics/UoN_gradient_logo_CMYK}
\end{figure}
\begin{figure}
\includegraphics[height=0.08\textheight]{../Graphics/UKRI_BBSR_Council-Logo_Horiz-CMYK}
\end{figure}
\end{multicols}
\makeatletter
Email: pcyra2@nottingham.ac.uk \\
Twitter: @AmoryResearch
\makeatother
\end{columns}
\end{frame}